\chapter{Reproducibilidad}
\label{ap:reproducibilidad}
% Modelo: anexo «Código fuente del proyecto» de la tesina de referencia (un enlace y la organización del repositorio),
% ampliado con los scripts de cada etapa y las semillas. Fuentes: Codigo/docs/USAGE.md, Codigo/CLAUDE.md.

El código del trabajo y los archivos \LaTeX{} de este documento están en el repositorio
\url{https://github.com/Gaviola/Tesina}. Este apéndice describe cómo está organizado y qué \emph{scripts} corresponden
a cada etapa de la experimentación. La guía de uso del repositorio, \texttt{Codigo/docs/USAGE.md}, detalla la
instalación del entorno, la descarga de los conjuntos de datos y los comandos de cada etapa.

\section{Organización del repositorio}
\label{sec:ap_repositorio}

El código vive en la carpeta \texttt{Codigo/}, organizada en cinco partes:
\begin{itemize}
    \item \textbf{\texttt{src/}}, el paquete con la carga de los datos, los modelos, el entrenamiento y las métricas
    de evaluación, que usan todos los \emph{scripts}.
    \item \textbf{\texttt{configs/}}, los archivos YAML con los hiperparámetros de cada modelo y de cada etapa, que
    son la única fuente de esos valores: los \emph{scripts} los leen y no los fijan en el código.
    \item \textbf{\texttt{scripts/}}, los programas que se ejecutan para entrenar, evaluar y producir las tablas y
    figuras de los Capítulos~\ref{cp:res_clasif} y~\ref{cp:res_gen}.
    \item \textbf{\texttt{slurm/}}, los trabajos que se envían a los clústeres TOKO y Mendieta
    (Sección~\ref{sec:met_entorno}).
    \item \textbf{\texttt{requirements.txt}} y \textbf{\texttt{scripts/setup\_env.py}}, que instalan las dependencias
    y eligen la versión de PyTorch según la GPU disponible.
\end{itemize}
Los conjuntos de datos y los resultados no se incluyen en el repositorio. Los conjuntos se descargan de sus fuentes
originales (Sección~\ref{sec:met_datasets}) a \texttt{Codigo/data/raw/}, y cada ejecución escribe sus resultados en
\texttt{Codigo/outputs/}.

\section{Scripts de cada etapa}
\label{sec:ap_scripts}

La Tabla~\ref{tab:ap_scripts} relaciona cada etapa de la experimentación con los \emph{scripts} principales que la
ejecutan, todos en \texttt{Codigo/scripts/}.

\begin{table}[H]
    \centering
    \caption{Scripts principales de cada etapa de la experimentación.}
    \label{tab:ap_scripts}
    \small
    \begin{tabularx}{\textwidth}{@{}>{\raggedright\arraybackslash}p{4.6cm}>{\raggedright\arraybackslash}X@{}}
        \toprule
        \textbf{Etapa} & \textbf{Scripts} \\
        \midrule
        Preparación de los datos y partición por lesión & \texttt{prepare\_data.py},
            \texttt{make\_lesion\_disjoint\_split.py} \\
        \filasep
        Clasificadores supervisados & \texttt{toko/submit\_sweep.py}, \texttt{train\_classifier.py},
            \texttt{lesion\_split\_report.py}, \texttt{classifier\_ensemble.py}, \texttt{melanoma\_detection.py} \\
        \filasep
        \glspl{vlm} en clasificación & \texttt{train\_vlm\_classify.py}, \texttt{eval\_dermlip\_classify.py},
            \texttt{eval\_dermlip\_linearprobe.py}, \texttt{dermlip\_split\_eval.py} \\
        \filasep
        Evaluación de la clasificación & \texttt{evaluate\_classify.py}, \texttt{statistical\_tests.py},
            \texttt{leakage\_analysis.py} \\
        \filasep
        Evaluaciones complementarias & \texttt{hiba\_por\_modalidad.py}, \texttt{robustness\_perturbations.py},
            \texttt{evaluate\_calibration.py}, \texttt{interpretability.py} \\
        \filasep
        Ajuste de los generadores & \texttt{train\_vlm\_stage1.py}, \texttt{train\_vlm\_stage2.py} \\
        \filasep
        \emph{Pipeline} y generación de reportes & \texttt{make\_dx\_csv.py}, \texttt{dermlip\_split\_eval.py},
            \texttt{build\_uncertainty\_3tier.py}, \texttt{generate\_report.py} \\
        \filasep
        Evaluación de los reportes & \texttt{evaluate\_reports.py}, \texttt{ensemble\_bootstrap.py},
            \texttt{uncertainty\_analysis.py}, \texttt{error\_taxonomy.py} \\
        \filasep
        Costo computacional & \texttt{measure\_cost.py}, \texttt{measure\_vlm\_memory.py} \\
        \bottomrule
    \end{tabularx}
\end{table}

\section{Semillas}
\label{sec:ap_semillas}

Todos los entrenamientos, la partición por lesión y la generación de los reportes fijan la semilla 42. Esa es también
la semilla de la ejecución de referencia de cada clasificador supervisado (Sección~\ref{sec:met_estadistica}), y las
otras dos ejecuciones (Sección~\ref{sec:met_supervisados_ensamble}) usan las semillas 43 y 44. Como la generación usa
decodificación voraz, la semilla no cambia los reportes, y se fija solo para que el entorno sea el mismo en todas las
ejecuciones.
